%!TEX root = ../main.tex

\begin{changemargin}{0.8cm}{0.8cm}

~\vfill{}

\section*{Abstrakt}
\vskip 0.5em

\sloppy
Tato práce představuje návrh, implementaci a vyhodnocení prototypu časové synchronizace založeného na mikrokontroléru pro budoucí experimenty s roji UAV. Vyvinutý systém kombinuje příjem absolutního času pomocí signálu DCF77, generování sekund podle RTK PPS, firmware pro STM32, diagnostickou komunikaci přes UART a záznam a analýzu dat v prostředí ROS2.

Hardwarová implementace je založena na vlastní desce plošných spojů s mikrokontrolérem STM32F042F6P6, napájením přes USB-C, 3.3 V regulací, vstupy pro DCF77 a RTK PPS, rozhraním SWD pro ladění a diagnostickým výstupem přes UART. Nativní USB komunikace mikrokontroléru STM32 byla během vývoje testována, ale ve finální ověřené verzi byl použit externí převodník CP2102 USB-UART, protože poskytoval spolehlivější komunikaci.

Firmware dekóduje a ověřuje rámce DCF77, udržuje softwarové hodiny reálného času, zpracovává přerušení RTK PPS a odesílá diagnostické zprávy do prostředí companion computeru. Uzly ROS2 přijímají UART zprávy, publikují je do témat, ukládají CSV logy a umožňují následné zpracování pomocí Python skriptů.

Finální vyhodnocení potvrdilo stabilní získávání času z DCF77 a generování sekund řízené signálem RTK PPS. Ve finálním RTK PPS testu zůstala měřená hodnota \texttt{RTK\_DIFF} rovna 0 ms pro všechny synchronizované vzorky při milisekundovém rozlišení měření a odhadovaný drift byl 0.000 ms/s. Výsledky neprokazují submikrosekundovou přesnost synchronizace, ale ověřují funkční synchronizační prototyp bez pozorovatelného dlouhodobého driftu softwarových hodin v dostupném rozlišení měření.
\vskip 1em

{\bf Klíčová slova:} \KlicovaSlova

\vskip 2.5cm

\end{changemargin}
